\chapter{Glosarium Istilah AI (A-Z)}
\label{chap:glossary}

Berikut adalah daftar istilah lengkap yang sering digunakan dalam dunia Artificial Intelligence, Machine Learning, dan Deep Learning.

\section*{A}
\begin{description}
    \item[A/B Testing] Metode eksperimen untuk membandingkan dua versi model atau sistem (versi A dan B) guna menentukan mana yang berkinerja lebih baik berdasarkan metrik tertentu.
    \item[Accuracy] Metrik evaluasi klasifikasi yang mengukur rasio prediksi benar terhadap total prediksi. Tidak disarankan untuk dataset yang tidak seimbang (imbalanced).
    \item[Activation Function] Fungsi matematis dalam neuron buatan yang menentukan apakah neuron tersebut harus "aktif" (mengirim sinyal) atau tidak. Contoh: ReLU, Sigmoid, Tanh.
    \item[AdaBoost (Adaptive Boosting)] Algoritma boosting yang melatih serangkaian weak learners (biasanya decision stumps) secara berurutan, di mana setiap model berfokus pada kesalahan model sebelumnya.
    \item[Adam (Adaptive Moment Estimation)] Algoritma optimasi populer yang menggabungkan ide Momentum dan RMSProp. Sangat efektif untuk deep learning.
    \item[Agent] Entitas dalam Reinforcement Learning yang berinteraksi dengan lingkungan (environment) untuk mencapai tujuan tertentu melalui tindakan (action).
    \item[AGI (Artificial General Intelligence)] Tingkat kecerdasan buatan hipotetis yang memiliki kemampuan kognitif setara atau melebihi manusia dalam berbagai tugas, bukan hanya satu tugas spesifik.
    \item[Algorithm] Serangkaian instruksi atau aturan logis yang diikuti komputer untuk memecahkan masalah atau melakukan perhitungan.
    \item[Alignment] Proses menyelaraskan perilaku sistem AI dengan nilai, etika, dan tujuan manusia.
    \item[Annotation] Proses memberikan label atau metadata pada data mentah (seperti gambar atau teks) agar dapat digunakan untuk melatih model supervised learning.
    \item[API (Application Programming Interface)] Antarmuka yang memungkinkan dua aplikasi perangkat lunak untuk berkomunikasi satu sama lain. Contoh: OpenAI API.
    \item[Architecture] Struktur desain model neural network, termasuk jumlah layer, jenis layer, dan koneksi antar neuron.
    \item[Artificial Intelligence (AI)] Cabang ilmu komputer yang berfokus pada penciptaan sistem yang mampu melakukan tugas yang biasanya memerlukan kecerdasan manusia.
    \item[Attention Mechanism] Teknik dalam deep learning (khususnya NLP) yang memungkinkan model untuk fokus pada bagian tertentu dari input saat menghasilkan output.
    \item[Autoencoder] Jenis neural network unsupervised yang digunakan untuk mempelajari representasi efisien (encoding) dari data input, biasanya untuk reduksi dimensi atau denoising.
    \item[AutoML (Automated Machine Learning)] Proses mengotomatisasi tugas-tugas end-to-end dalam penerapan machine learning, mulai dari preprocessing hingga pemilihan model.
    \item[Autoregressive Model] Model statistik yang memprediksi nilai masa depan berdasarkan nilai masa lalu. Dalam LLM, ini berarti memprediksi kata berikutnya berdasarkan kata-kata sebelumnya.
\end{description}

\section*{B}
\begin{description}
    \item[Backpropagation] Algoritma utama untuk melatih neural network. Menghitung gradien loss function terhadap bobot jaringan dengan menggunakan aturan rantai (chain rule), lalu memperbarui bobot untuk meminimalkan error.
    \item[Bagging (Bootstrap Aggregating)] Teknik ensemble learning yang melatih beberapa model secara independen pada subset data acak (bootstrap) dan menggabungkan prediksinya (biasanya dirata-rata). Contoh: Random Forest.
    \item[Batch] Sekumpulan sampel data yang diproses oleh model dalam satu iterasi pelatihan sebelum memperbarui bobot.
    \item[Batch Normalization] Teknik untuk menstabilkan dan mempercepat pelatihan deep neural network dengan menormalkan input layer untuk setiap mini-batch.
    \item[Bayes' Theorem] Teorema probabilitas yang menggambarkan probabilitas suatu peristiwa berdasarkan pengetahuan sebelumnya tentang kondisi yang mungkin terkait.
    \item[Bayesian Network] Model grafis probabilistik yang merepresentasikan sekumpulan variabel dan dependensi kondisionalnya melalui Directed Acyclic Graph (DAG).
    \item[Beam Search] Algoritma pencarian heuristik yang mengeksplorasi graf dengan memperluas node yang paling menjanjikan dalam jumlah terbatas (beam width). Digunakan dalam decoding model bahasa.
    \item[Bias (Inductive Bias)] Asumsi yang dibuat oleh model untuk menyederhanakan pembelajaran fungsi target.
    \item[Bias (Ethical)] Prasangka sistematis dalam hasil model AI yang mencerminkan bias manusia atau data historis.
    \item[Bias (Neuron)] Parameter tambahan dalam neuron (selain bobot) yang memungkinkan fungsi aktivasi digeser ke kiri atau kanan.
    \item[Big Data] Kumpulan data yang sangat besar, kompleks, dan berkembang cepat yang sulit diproses menggunakan metode tradisional.
    \item[Binary Classification] Tugas klasifikasi dengan hanya dua kemungkinan kelas output (misal: Spam/Bukan Spam).
    \item[Black Box] Sistem AI yang proses internal atau pengambilan keputusannya sulit atau tidak mungkin dipahami oleh manusia.
    \item[Boosting] Teknik ensemble learning yang menggabungkan beberapa weak learners secara berurutan untuk membentuk strong learner.
    \item[Bounding Box] Kotak persegi panjang yang digunakan dalam deteksi objek untuk menandai lokasi objek dalam gambar.
\end{description}

\section*{C}
\begin{description}
    \item[Capsule Network (CapsNet)] Jenis neural network yang dirancang untuk mengatasi kekurangan CNN dalam menangani hierarki spasial dan pose objek.
    \item[Categorical Data] Data yang merepresentasikan kategori atau label (misal: Warna = Merah, Biru, Hijau).
    \item[Chain of Thought (CoT)] Teknik prompting di mana model didorong untuk menghasilkan langkah-langkah penalaran perantara sebelum memberikan jawaban akhir.
    \item[Chatbot] Program komputer yang dirancang untuk mensimulasikan percakapan dengan pengguna manusia, seringkali menggunakan NLP.
    \item[Classification] Tugas supervised learning di mana model memprediksi label kategori diskrit untuk data input.
    \item[Clustering] Tugas unsupervised learning untuk mengelompokkan data yang mirip ke dalam cluster tanpa label sebelumnya.
    \item[CNN (Convolutional Neural Network)] Jenis neural network yang sangat efektif untuk memproses data grid seperti gambar, menggunakan operasi konvolusi.
    \item[Computer Vision] Bidang AI yang melatih komputer untuk menafsirkan dan memahami dunia visual dari gambar atau video.
    \item[Confusion Matrix] Tabel yang digunakan untuk mengevaluasi kinerja model klasifikasi, menunjukkan True Positives, False Positives, True Negatives, dan False Negatives.
    \item[Context Window] Jumlah maksimum token yang dapat diproses oleh model bahasa sekaligus (sebagai input prompt).
    \item[Continuous Learning] Kemampuan sistem AI untuk terus belajar dari data baru tanpa melupakan pengetahuan lama (menghindari catastrophic forgetting).
    \item[Corpus] Kumpulan besar teks tertulis atau lisan yang digunakan untuk melatih model bahasa.
    \item[Cost Function] Sinonim untuk Loss Function; fungsi yang mengukur seberapa buruk kinerja model.
    \item[Cross-Validation] Teknik statistik untuk menilai kinerja model dengan membagi data menjadi beberapa subset (folds) dan melatih/menguji secara bergantian.
\end{description}

\section*{D}
\begin{description}
    \item[Data Augmentation] Teknik meningkatkan keragaman data latih dengan menerapkan transformasi acak (seperti rotasi, cropping, flipping) pada data yang ada.
    \item[Data Leakage] Masalah di mana informasi dari data uji bocor ke dalam proses pelatihan, menyebabkan estimasi kinerja yang terlalu optimis.
    \item[Data Mining] Proses menemukan pola, korelasi, atau anomali dalam dataset besar.
    \item[Data Science] Bidang interdisipliner yang menggunakan metode ilmiah, proses, algoritma, dan sistem untuk mengekstrak pengetahuan dari data.
    \item[Dataset] Kumpulan data yang terorganisir.
    \item[Decision Boundary] Permukaan pemisah yang membedakan kelas yang berbeda dalam ruang fitur.
    \item[Decision Tree] Model prediksi berbentuk struktur pohon di mana setiap node internal merepresentasikan tes pada atribut, setiap cabang merepresentasikan hasil tes, dan setiap daun merepresentasikan label kelas.
    \item[Decoder] Bagian dari neural network (seperti Transformer atau Autoencoder) yang mengubah representasi internal kembali menjadi output data (misal: teks atau gambar).
    \item[Deep Learning] Subbidang machine learning yang menggunakan neural network dengan banyak layer (deep) untuk mempelajari representasi data yang kompleks.
    \item[Deep Q-Network (DQN)] Algoritma reinforcement learning yang menggabungkan Q-Learning dengan deep neural networks.
    \item[Deepfake] Media sintetis (biasanya video atau audio) di mana kemiripan seseorang diganti dengan orang lain menggunakan deep learning.
    \item[Dense Layer] Layer dalam neural network di mana setiap neuron terhubung ke setiap neuron di layer sebelumnya (fully connected).
    \item[Diffusion Model] Model generatif yang belajar menghasilkan data dengan membalikkan proses difusi noise (menambahkan noise bertahap hingga data hancur, lalu belajar merekonstruksinya).
    \item[Dimensionality Reduction] Proses mengurangi jumlah variabel acak atau fitur dalam dataset, misal menggunakan PCA atau t-SNE.
    \item[Discriminator] Dalam GAN, ini adalah neural network yang bertugas membedakan antara data asli dan data palsu yang dibuat oleh Generator.
    \item[Dropout] Teknik regularisasi di mana neuron dipilih secara acak untuk diabaikan (dropped out) selama pelatihan untuk mencegah overfitting.
\end{description}

\section*{E}
\begin{description}
    \item[Early Stopping] Teknik regularisasi yang menghentikan pelatihan model ketika kinerja pada validation set mulai menurun, meskipun training loss masih turun.
    \item[Embedding] Representasi vektor padat (dense vector) dari data diskrit (seperti kata atau item) dalam ruang kontinu berdimensi rendah, di mana item serupa berdekatan.
    \item[Encoder] Bagian dari neural network yang memproses input mentah menjadi representasi internal (embedding) atau fitur.
    \item[Ensemble Learning] Metode yang menggabungkan prediksi dari beberapa model untuk meningkatkan kinerja keseluruhan.
    \item[Entropy] Ukuran ketidakteraturan atau ketidakpastian dalam sistem. Dalam teori informasi, mengukur jumlah informasi rata-rata.
    \item[Epoch] Satu putaran penuh melalui seluruh dataset pelatihan.
    \item[Expert System] Sistem komputer yang meniru kemampuan pengambilan keputusan seorang ahli manusia, biasanya menggunakan aturan if-then.
    \item[Explainable AI (XAI)] Metode dan teknik yang memungkinkan hasil solusi AI dipahami oleh manusia.
\end{description}

\section*{F}
\begin{description}
    \item[F1 Score] Harmonic mean dari Precision dan Recall. Berguna untuk dataset yang tidak seimbang.
    \item[False Negative] Kesalahan di mana model memprediksi kelas negatif padahal sebenarnya positif.
    \item[False Positive] Kesalahan di mana model memprediksi kelas positif padahal sebenarnya negatif.
    \item[Feature] Atribut atau properti individual dari data yang diamati (misal: tinggi badan, warna piksel).
    \item[Feature Engineering] Proses menggunakan pengetahuan domain untuk mengekstrak fitur yang relevan dari data mentah agar algoritma ML bekerja lebih baik.
    \item[Feature Extraction] Proses otomatis atau manual untuk mereduksi dimensi data mentah menjadi fitur yang lebih mudah dikelola.
    \item[Feature Selection] Proses memilih subset fitur yang paling relevan untuk digunakan dalam pembangunan model.
    \item[Federated Learning] Teknik melatih algoritma ML di beberapa perangkat lokal terdesentralisasi tanpa menukar sampel data antar perangkat.
    \item[Feedforward Neural Network] Jenis neural network paling sederhana di mana koneksi antar node tidak membentuk siklus.
    \item[Few-Shot Learning] Kemampuan model untuk belajar tugas baru dengan hanya melihat sedikit contoh (sedikit shot).
    \item[Fine-Tuning] Proses mengambil model yang sudah dilatih (pre-trained) dan melatihnya lebih lanjut pada dataset spesifik untuk tugas tertentu.
    \item[Foundation Model] Model AI skala besar (seperti GPT-4, BERT) yang dilatih pada data yang luas dan dapat diadaptasi untuk berbagai tugas hilir.
\end{description}

\section*{G}
\begin{description}
    \item[GAN (Generative Adversarial Network)] Arsitektur deep learning yang terdiri dari dua jaringan (Generator dan Discriminator) yang bersaing satu sama lain.
    \item[Gated Recurrent Unit (GRU)] Varian dari LSTM yang lebih sederhana namun seringkali sama efektifnya.
    \item[Generalization] Kemampuan model untuk berkinerja baik pada data baru yang tidak pernah dilihat sebelumnya (data uji).
    \item[Generative AI] Cabang AI yang berfokus pada pembuatan konten baru (teks, gambar, audio) yang mirip dengan data latih.
    \item[Genetic Algorithm] Algoritma pencarian heuristik yang terinspirasi oleh proses seleksi alam dan genetika.
    \item[Gradient] Vektor turunan parsial yang menunjukkan arah kenaikan tercepat dari suatu fungsi.
    \item[Gradient Descent] Algoritma optimasi iteratif untuk menemukan minimum lokal dari fungsi yang dapat didiferensiasi (biasanya loss function).
    \item[Graph Neural Network (GNN)] Neural network yang dirancang untuk bekerja langsung pada struktur data graf.
    \item[Ground Truth] Informasi yang diketahui sebagai kebenaran mutlak, digunakan untuk melatih dan menguji model (label yang benar).
\end{description}

\section*{H}
\begin{description}
    \item[Hallucination] Fenomena di mana model bahasa menghasilkan informasi yang tampak masuk akal tetapi sebenarnya salah atau tidak berdasar fakta.
    \item[Heuristic] Aturan praktis atau metode jalan pintas yang digunakan untuk memecahkan masalah lebih cepat, meskipun tidak menjamin solusi optimal.
    \item[Hidden Layer] Layer neuron yang berada di antara input layer dan output layer dalam neural network.
    \item[Human-in-the-Loop (HITL)] Model interaksi di mana manusia membantu melatih, meninjau, atau mengoreksi keputusan sistem AI.
    \item[Hyperparameter] Parameter konfigurasi model yang ditetapkan sebelum pelatihan dimulai (misal: learning rate, jumlah layer) dan tidak dipelajari dari data.
    \item[Hyperparameter Tuning] Proses mencari kombinasi hyperparameter terbaik untuk memaksimalkan kinerja model.
\end{description}

\section*{I}
\begin{description}
    \item[Image Segmentation] Tugas mempartisi gambar digital menjadi beberapa segmen (kumpulan piksel) untuk menyederhanakan representasi gambar.
    \item[Inference] Proses menggunakan model yang sudah dilatih untuk membuat prediksi pada data baru.
    \item[Instance Segmentation] Jenis segmentasi gambar yang membedakan setiap instance objek individu (misal: membedakan antara dua mobil yang berbeda).
    \item[IoU (Intersection over Union)] Metrik evaluasi untuk deteksi objek yang mengukur tumpang tindih antara kotak prediksi dan ground truth.
\end{description}

\section*{J}
\begin{description}
    \item[Jaccard Index] Ukuran statistik untuk kemiripan antara set sampel (mirip dengan IoU).
    \item[Joint Probability] Probabilitas dua peristiwa terjadi secara bersamaan.
    \item[Jupyter Notebook] Aplikasi web open-source yang memungkinkan pembuatan dan pembagian dokumen yang berisi kode langsung, persamaan, visualisasi, dan teks naratif.
\end{description}

\section*{K}
\begin{description}
    \item[K-Fold Cross-Validation] Teknik validasi di mana data dibagi menjadi K subset; model dilatih K kali, setiap kali menggunakan satu subset berbeda sebagai data uji.
    \item[K-Means Clustering] Algoritma clustering populer yang mempartisi data menjadi K cluster berdasarkan jarak rata-rata ke centroid.
    \item[K-Nearest Neighbors (KNN)] Algoritma klasifikasi sederhana yang mengklasifikasikan sampel data berdasarkan mayoritas label dari K tetangga terdekatnya.
    \item[Kernel Trick] Metode yang memungkinkan algoritma linear (seperti SVM) untuk memecahkan masalah non-linear dengan memetakan input ke ruang dimensi yang lebih tinggi.
    \item[Knowledge Distillation] Proses mentransfer pengetahuan dari model besar (teacher) ke model kecil (student) agar model kecil bisa berkinerja mendekati model besar namun lebih efisien.
    \item[Knowledge Graph] Representasi terstruktur dari fakta yang terdiri dari entitas, hubungan, dan informasi semantik.
\end{description}

\section*{L}
\begin{description}
    \item[Label] Output yang diinginkan atau jawaban yang benar untuk suatu sampel data dalam supervised learning.
    \item[Latent Space] Ruang representasi abstrak (biasanya berdimensi lebih rendah) di mana data input dipetakan oleh model (seperti encoder).
    \item[Layer] Kumpulan neuron yang memproses fitur input secara bersamaan.
    \item[Learning Rate] Hyperparameter yang mengontrol seberapa besar langkah yang diambil model saat memperbarui bobot selama pelatihan.
    \item[Linear Regression] Algoritma statistik untuk memodelkan hubungan antara variabel dependen skalar dan satu atau lebih variabel independen.
    \item[LLM (Large Language Model)] Model bahasa dengan parameter miliaran yang dilatih pada dataset teks yang sangat besar.
    \item[Logistic Regression] Algoritma klasifikasi statistik yang digunakan untuk memprediksi probabilitas kejadian biner.
    \item[Loss Function] Fungsi yang menghitung selisih antara prediksi model dan label sebenarnya. Tujuannya adalah meminimalkan fungsi ini.
    \item[LSTM (Long Short-Term Memory)] Jenis RNN khusus yang mampu mempelajari ketergantungan jangka panjang, mengatasi masalah vanishing gradient.
\end{description}

\section*{M}
\begin{description}
    \item[Machine Learning (ML)] Subbidang AI yang memberi komputer kemampuan untuk belajar tanpa diprogram secara eksplisit.
    \item[Machine Translation] Penggunaan software untuk menerjemahkan teks atau ucapan dari satu bahasa ke bahasa lain.
    \item[Masked Language Modeling (MLM)] Teknik pelatihan pre-training (seperti di BERT) di mana beberapa kata dalam kalimat disembunyikan (masked) dan model harus memprediksinya.
    \item[Mean Squared Error (MSE)] Loss function standar untuk regresi, menghitung rata-rata kuadrat selisih antara nilai prediksi dan nilai sebenarnya.
    \item[Meta-Learning] "Learning to learn"; pendekatan di mana model belajar cara mempelajari tugas baru dengan cepat.
    \item[Metric] Ukuran kuantitatif untuk menilai kinerja model (misal: akurasi, presisi, recall).
    \item[MLOps (Machine Learning Operations)] Praktik untuk menyatukan pengembangan sistem ML (Dev) dan operasi sistem ML (Ops).
    \item[Model] Representasi matematis dari pola yang dipelajari dari data.
    \item[Model Drift] Penurunan kinerja model seiring waktu karena perubahan distribusi data di dunia nyata.
    \item[Momentum] Teknik optimasi yang mempercepat Gradient Descent dengan menambahkan fraksi dari vektor update sebelumnya ke update saat ini.
    \item[Multi-Head Attention] Komponen Transformer yang memungkinkan model memperhatikan informasi dari berbagai representasi subruang secara bersamaan.
    \item[Multimodal] Sistem AI yang dapat memproses dan menghubungkan informasi dari berbagai modalitas (teks, gambar, audio).
\end{description}

\section*{N}
\begin{description}
    \item[Naive Bayes] Algoritma klasifikasi probabilistik berdasarkan Teorema Bayes dengan asumsi independensi yang kuat (naif) antar fitur.
    \item[Natural Language Processing (NLP)] Cabang AI yang berfokus pada interaksi antara komputer dan bahasa manusia.
    \item[Neural Network] Model komputasi yang terinspirasi oleh struktur otak biologis, terdiri dari neuron yang saling terhubung.
    \item[Normalization] Proses penskalaan nilai input (biasanya ke rentang 0-1) agar memiliki rentang yang serupa.
\end{description}

\section*{O}
\begin{description}
    \item[Object Detection] Tugas CV untuk mengidentifikasi dan melokalisasi (dengan bounding box) objek dalam gambar.
    \item[One-Hot Encoding] Representasi variabel kategorikal sebagai vektor biner di mana hanya satu elemen yang bernilai 1 (hot).
    \item[Open Source] Software yang kode sumbernya tersedia untuk umum untuk digunakan, dimodifikasi, dan didistribusikan.
    \item[Optimization] Proses menyesuaikan parameter model untuk meminimalkan loss function.
    \item[Outlier] Titik data yang berbeda secara signifikan dari pengamatan lainnya.
    \item[Overfitting] Kondisi di mana model belajar detail dan noise pada data latih terlalu baik sehingga kinerjanya buruk pada data baru.
\end{description}

\section*{P}
\begin{description}
    \item[Parameter] Variabel internal model yang dipelajari dari data (misal: bobot dan bias).
    \item[Perceptron] Algoritma neural network paling sederhana (satu neuron) untuk klasifikasi biner linear.
    \item[Pooling] Operasi dalam CNN untuk mengurangi dimensi spasial representasi (misal: Max Pooling).
    \item[Precision] Rasio prediksi positif yang benar terhadap total prediksi positif. "Seberapa akurat saat model bilang 'Ya'?"
    \item[Pre-training] Tahap awal pelatihan model pada dataset besar umum sebelum di-fine-tune untuk tugas khusus.
    \item[Principal Component Analysis (PCA)] Teknik reduksi dimensi yang mengubah fitur berkorelasi menjadi sekumpulan fitur ortogonal (komponen utama).
    \item[Prompt Engineering] Seni merancang input (prompt) untuk mendapatkan output terbaik dari model bahasa generatif.
    \item[Pruning] Teknik mengompresi neural network dengan menghapus bobot atau neuron yang kurang penting.
    \item[PyTorch] Library open-source untuk machine learning yang dikembangkan oleh Facebook (Meta), populer karena fleksibilitasnya.
\end{description}

\section*{Q}
\begin{description}
    \item[Q-Learning] Algoritma reinforcement learning off-policy yang bertujuan mempelajari fungsi nilai aksi (Q-function).
    \item[Quantization] Teknik mengurangi presisi representasi angka (misal: dari 32-bit float ke 8-bit integer) untuk memperkecil ukuran model dan mempercepat inferensi.
\end{description}

\section*{R}
\begin{description}
    \item[RAG (Retrieval-Augmented Generation)] Teknik meningkatkan akurasi LLM dengan mengambil data relevan dari sumber eksternal sebelum menghasilkan jawaban.
    \item[Random Forest] Algoritma ensemble yang terdiri dari banyak decision trees.
    \item[Recall (Sensitivity)] Rasio prediksi positif yang benar terhadap total positif yang sebenarnya. "Seberapa banyak 'Ya' yang berhasil ditemukan?"
    \item[Recommendation System] Sistem yang memprediksi preferensi pengguna terhadap item (misal: film, produk).
    \item[Rectified Linear Unit (ReLU)] Fungsi aktivasi non-linear yang outputnya adalah input jika positif, dan 0 jika negatif. $f(x) = \max(0, x)$.
    \item[Recurrent Neural Network (RNN)] Neural network dengan loop yang memungkinkan informasi bertahan, cocok untuk data sekuensial.
    \item[Regression] Tugas supervised learning untuk memprediksi nilai kontinu (angka).
    \item[Regularization] Teknik untuk mencegah overfitting dengan menambahkan penalti pada loss function (misal: L1, L2).
    \item[Reinforcement Learning (RL)] Jenis ML di mana agen belajar membuat keputusan dengan melakukan tindakan dan menerima reward atau punishment.
    \item[RLHF (Reinforcement Learning from Human Feedback)] Metode melatih model bahasa dengan menggunakan feedback manusia sebagai sinyal reward.
    \item[ROC Curve (Receiver Operating Characteristic)] Grafik yang menggambarkan kinerja model klasifikasi pada berbagai ambang batas.
\end{description}

\section*{S}
\begin{description}
    \item[Self-Attention] Mekanisme yang memungkinkan model menghubungkan posisi yang berbeda dari satu urutan untuk menghitung representasi urutan tersebut.
    \item[Self-Supervised Learning] Metode pembelajaran di mana data menyediakan labelnya sendiri (misal: memprediksi bagian gambar yang hilang).
    \item[Semantic Segmentation] Tugas CV yang mengklasifikasikan setiap piksel dalam gambar ke dalam kelas tertentu.
    \item[Sentiment Analysis] Penggunaan NLP untuk mengidentifikasi dan mengekstrak opini atau sentimen dari teks.
    \item[Sequence-to-Sequence (Seq2Seq)] Model yang mengubah urutan input menjadi urutan output (misal: terjemahan bahasa).
    \item[Sigmoid Function] Fungsi aktivasi berbentuk S yang memetakan input ke rentang (0, 1).
    \item[Softmax Function] Fungsi yang mengubah vektor angka menjadi vektor probabilitas yang jumlahnya 1.
    \item[Sparse Matrix] Matriks di mana sebagian besar elemennya adalah nol.
    \item[Stochastic Gradient Descent (SGD)] Varian gradient descent yang memperbarui bobot berdasarkan gradien dari satu sampel acak (atau mini-batch) per langkah.
    \item[Stride] Langkah pergeseran filter konvolusi melintasi input.
    \item[Supervised Learning] Paradigma ML di mana model dilatih pada data berlabel (input-output pairs).
    \item[Support Vector Machine (SVM)] Algoritma supervised learning yang mencari hyperplane terbaik yang memisahkan kelas dengan margin terbesar.
    \item[Synthetic Data] Data yang dibuat secara artifisial, bukan diperoleh dari pengukuran langsung.
\end{description}

\section*{T}
\begin{description}
    \item[Tanh (Hyperbolic Tangent)] Fungsi aktivasi yang memetakan input ke rentang (-1, 1).
    \item[Tensor] Objek matematika yang menggeneralisasi skalar, vektor, dan matriks ke dimensi yang lebih tinggi.
    \item[TensorFlow] Library ML open-source yang dikembangkan oleh Google.
    \item[Test Set] Subset data yang tidak pernah dilihat model selama pelatihan, digunakan untuk evaluasi akhir.
    \item[Token] Unit dasar teks (bisa berupa kata, subkata, atau karakter) yang diproses oleh model bahasa.
    \item[Tokenization] Proses memecah teks menjadi token-token individual.
    \item[Training Set] Subset data yang digunakan untuk melatih model.
    \item[Transfer Learning] Teknik menggunakan model yang sudah dilatih pada satu tugas sebagai titik awal untuk tugas lain yang terkait.
    \item[Transformer] Arsitektur deep learning yang mengandalkan mekanisme self-attention, menjadi dasar bagi sebagian besar model NLP modern.
    \item[True Negative] Prediksi negatif yang benar.
    \item[True Positive] Prediksi positif yang benar.
    \item[Turing Test] Tes kemampuan mesin untuk menunjukkan perilaku cerdas yang tidak dapat dibedakan dari manusia.
\end{description}

\section*{U}
\begin{description}
    \item[Underfitting] Kondisi di mana model terlalu sederhana untuk menangkap pola dasar data (kinerja buruk pada training dan test set).
    \item[Unsupervised Learning] Paradigma ML di mana model mencari pola dalam data yang tidak berlabel.
\end{description}

\section*{V}
\begin{description}
    \item[Validation Set] Subset data yang digunakan untuk menyetel hyperparameter dan mencegah overfitting selama pelatihan.
    \item[Vanishing Gradient Problem] Masalah dalam melatih deep neural network di mana gradien menjadi sangat kecil, sehingga bobot di layer awal berhenti berubah.
    \item[Variance] Ukuran seberapa jauh prediksi model tersebar dari nilai rata-ratanya. Variance tinggi sering dikaitkan dengan overfitting.
    \item[Vector] Array satu dimensi angka.
    \item[Vision Transformer (ViT)] Adaptasi arsitektur Transformer untuk tugas visi komputer.
\end{description}

\section*{W}
\begin{description}
    \item[Weight] Parameter dalam neural network yang mengubah data input di dalam hidden layers. Kekuatan koneksi antar neuron.
    \item[Word Embedding] Representasi vektor kata di mana kata-kata dengan makna serupa memiliki representasi vektor yang dekat.
\end{description}

\section*{X}
\begin{description}
    \item[XGBoost] Implementasi gradient boosting yang efisien dan terukur, sangat populer dalam kompetisi data science (Kaggle).
\end{description}

\section*{Z}
\begin{description}
    \item[Zero-Shot Learning] Kemampuan model untuk melakukan tugas yang tidak pernah dilatihkan secara eksplisit.
\end{description}

\section*{Simbol Matematika Umum}
\begin{itemize}
    \item $\sum$ : Summation (Penjumlahan)
    \item $\prod$ : Product (Perkalian beruntun)
    \item $\nabla$ : Nabla (Gradien)
    \item $\partial$ : Turunan parsial
    \item $\mathbb{R}$ : Himpunan bilangan riil
    \item $\in$ : Anggota himpunan
    \item $\hat{y}$ : Nilai prediksi
    \item $\mu$ : Mean (Rata-rata populasi)
    \item $\sigma$ : Standar deviasi
    \item $\theta$ : Parameter model
    \item $\alpha$ : Learning rate
    \item $\lambda$ : Parameter regularisasi
\end{itemize}
